\documentclass[11pt]{article}

% Packages
\usepackage[top=0.75in, bottom=1in, left=0.625in, right=0.625in]{geometry}
\usepackage{newtxtext, newtxmath}
\usepackage{amsmath}
\usepackage{graphicx}
% \usepackage{minipage}
\usepackage{caption}
% \usepackage{float}
\usepackage{hyperref}
\usepackage{cite}
\usepackage{listings}
\usepackage{xcolor}

% Figure Path
\graphicspath{ {./figures/} }

% Code listing style
\lstset{
    basicstyle=\ttfamily\small,
    keywordstyle=\color{blue},
    commentstyle=\color{gray},
    stringstyle=\color{teal},
    breaklines=true,
    frame=single,
    numbers=left,
    numberstyle=\tiny\color{gray}
}

% Hyperlink style
\hypersetup{
    colorlinks=true,
    linkcolor=black,
    citecolor=black,
    urlcolor=blue
}

% -----------------------------------------------
%   Title block
% -----------------------------------------------
\title{MPI Parallelisation For Computing The Mandelbrot Set}
\author{
    Jacob Gallop \\
    \small Student ID: 20253130 \\
    % \small High Performance Computing \\
    % \small Curtin University
}
% \date{\today}

% -----------------------------------------------
\begin{document}

\maketitle

% - Describe mandelbrot set and algorithm
% - Include render
% - Explain serial changes
% - Show serial time

% Static Decomposition (Block Partitioning)
% - Describe data distribution
% - Explain algorithm for distributing work
% - Code snippets showing changes
% - Compare timing to serial code, and to ideal speedup (Amdahl's Law)
% - Timing figure. Explain

% Dynamic Decomposition (Master-worker)
% - Include previous points
% - [Non-blocking algorithm]
% - Figure of runtime vs. chunksize. Explain
% - Compare timing of chunksize /= 10 to static. Explain.

% Static Decomposition (Cyclic Partitioning)
% - Include previous points
% - Compare timing/load-balancing figure to block partitioning. Explain.

% Conclusion
% - Comparison of three methods implemented
% - Strengths and weaknesses of each.
% - Best?


\begin{abstract}
\vspace{2em}
\begin{minipage}{1.0\linewidth}
\vspace{-2em}
The Mandelbrot set is a famous fractal pattern produced by iterating the complex quadratic recurrence relation:
\begin{equation}\label{eq:set}
    Z_{n+1} = Z_{n}^{2} + C
\end{equation}
where $Z_{n}$ denotes the complex value after $n$ iterations, and $C$ 
denotes the complex point. After $n = N$ iterations, where $N$ is 
defined as an upper-bound for number of iterations, points which 
have not diverged form the Mandelbrot set.

It can be assumed that points which exceed the complex-plane origin 
by a radius of $2$ will eventually diverge, yielding the 
following axiom for computing the Mandelbrot set $M$:
\begin{equation}\label{eq:axiom}
    \mathbb{M} = \{\;c \in \mathbb{C},\;\sup_{n \in \mathbb{N}} |Z_{n}| \leq 2,\;\text{where}\;Z_{0}=0\;\text{and}\;Z_{n+1}=Z_{n}^{2}+c\;\}
\end{equation}

. Points of the Mandelbrot set will inevitably  $N$ 

It can be shown that values of $Z_{n}$ which exceed $2$ will diverge, as such 
The value of $Z_{0}$ is usually taken to be zero. The iteration is repeated until the magnitude of $Z_{n}$ ($|Z_{n}|$) exceeds $2$ (at which point the sequence diverges), or a maximum number of iterations is reached (at
which point the point is considered to be part of the Mandelbrot set). The iteration
count is then used to assign a colour to the pixel representing that point, thus
rendering the set.
\end{minipage}
\hfill

% \fbox
% {
% \begin{minipage}[c][0.6\linewidth][c]{0.6\linewidth}
\begin{figure}[h]
    \centering
    \includegraphics[width=0.5\linewidth]{mandelbrot.jpg}
    \captionof{figure}{Render of the Mandelbrot Set}
    \label{fig:mandelbrot}
\end{figure}    
% \end{minipage}
% }

\end{abstract}

% -----------------------------------------------
\section{Introduction}
\label{sec:intro}

While the scope of this paper focuses on the parallel implementations for computing 
Mandelbrot set points, the following optimisations should be acknowledged separately as 
they apply to the serial algorithm, and are therefore common to all MPI implementations. 

\textbf{Breaking up "complex" number types into parametric values} \\
Complex values are iterated in parametric form by the following algorithm:
\begin{lstlisting}[language=C, caption={Parametric Mandelbrot set Algorithm}]
for (loop = 0; loop < N*N; loop++) {
    i = loop % N;
    j = loop / N;

    ki = zi = (4.0f * (i - (float)N/2)) / N;
    kj = zj = (4.0f * (j - (float)N/2)) / N;

    k = 1;
    while (((zi * zi) + (zj * zj) <= 4.0f) && (k++ < MAXITER)) { 
        new_zi = (zi * zi) - (zj * zj) + ki;
        zj = 2.0f * zi * zj + kj;
        zi = new_zi;
    }
  
    x[loop] = logf((float)k) / logf((float)MAXITER);
}
\end{lstlisting}
Most notably, the loop conditional only compares the complex component's 
squared sum to $4$, as opposed to comparing \verb|cabs(z)| and computing the 
square root of the complex component's squared sum.

Additionally, replacing the complex type with real number types may yield faster 
arithmetic operations from the reduced overhead and function calls in the "hot-loop"; 
however, it is difficult to ascertain given the how effectively modern C compilers 
can optimise arithmetic. \\

\textbf{Mirroring about the real-axis}
Due to the symmetry of the Mandelbrot set about the real number axis, points for 
only one half must be iterated; the other half can simply mirror the results. This 
led to the memory usage to be roughly halved, and  

% -----------------------------------------------
\section{Background}
\label{sec:background}

Summarise related work and provide the theoretical or technical
context that the reader needs to understand your contribution.

% -----------------------------------------------
\section{Methodology}
\label{sec:methodology}

Describe your approach, design decisions, algorithms, or system
architecture. Use sub-sections as needed.

\subsection{System Design}

Details of the design here.

\subsection{Algorithm}

You can typeset algorithms with pseudo-code:

\begin{lstlisting}[language=Python, caption={Example algorithm}]
def example(n):
    result = 0
    for i in range(n):
        result += i   # O(n)
    return result
\end{lstlisting}

% -----------------------------------------------
\section{Evaluation}
\label{sec:evaluation}

Present your experimental setup, metrics, datasets, and results.
Tables and figures belong here.

\begin{table}[h]
    \centering
    \begin{tabular}{lcc}
        \hline
        \textbf{Method} & \textbf{Accuracy (\%)} & \textbf{Time (s)} \\
        \hline
        Baseline        & 72.4                   & 0.32              \\
        Proposed        & 89.1                   & 0.41              \\
        \hline
    \end{tabular}
    \caption{Comparison of methods.}
    \label{tab:results}
\end{table}

% -----------------------------------------------
\section{Discussion}
\label{sec:discussion}

Interpret your results, discuss limitations, and consider directions
for future work.

% -----------------------------------------------
\section{Conclusion}
\label{sec:conclusion}

Summarise the contributions and key findings of the paper.

% -----------------------------------------------
\bibliographystyle{plain}
\bibliography{references}   % references.bib

% If you prefer inline references, use thebibliography instead:
%
% \begin{thebibliography}{9}
% \bibitem{knuth1997}
%     D.~E. Knuth, \textit{The Art of Computer Programming}, Vol.~1.
%     Addison-Wesley, 1997.
% \end{thebibliography}

\end{document}

